% Web専用原稿．独立査読・所蔵文献該当箇所の照合は未了．
\section{分解体と代数閉包}
% genron:begin K-001406
% genron:use K-000001
% genron:use K-000002
% genron:use K-000003
% genron:use K-000004
% genron:use K-000005
% genron:use K-000007
% genron:use K-000017
% genron:use K-000024
% genron:use K-001190
% genron:use K-001195
% genron:use K-000356
\begin{defn}[分解体]
体 $K$ 上の非零多項式 $f$ が拡大体 $L$ で一次式の積に分解し，$L$ が $K$ と $f$ の根によって生成されるとき，$L$ を $f$ の $K$ 上の分解体という．（参照：\GenronRef{K-001190}，\GenronRef{K-001195}．）
\end{defn}
% genron:end K-001406
% genron:begin K-001407
% genron:use K-000001
% genron:use K-000002
% genron:use K-000003
% genron:use K-000004
% genron:use K-000005
% genron:use K-000007
% genron:use K-000017
% genron:use K-000024
% genron:use K-001406
% genron:use K-001237
% genron:use K-001238
% genron:use K-001310
% genron:use K-001191
% genron:use K-001190
\begin{prop}[分解体の存在]
任意の非零多項式 $f\in K[X]$ に対して有限次の分解体が存在する．（参照：\GenronRef{K-001190}，\GenronRef{K-001406}．）
\end{prop}
\begin{proof}
次数について帰納する．定数なら $K$ でよい．正次数なら既約因子 $p$ を選び，$K[X]/(p)$ において $p$ の根 $\alpha$ を得る．そこで $f=(X-\alpha)g$ と分解する．帰納法により $g$ の有限次分解体を得れば，その体は $f$ の全根を含む．その中で $K$ と全根から生成される部分体は求める分解体であり，有限次性は塔の公式から従う．
\end{proof}
% genron:end K-001407
% genron:begin K-001408
% genron:use K-000001
% genron:use K-000002
% genron:use K-000003
% genron:use K-000004
% genron:use K-000005
% genron:use K-000007
% genron:use K-000017
% genron:use K-000024
% genron:use K-001307
% genron:use K-001308
% genron:use K-001238
% genron:use K-000331
\begin{prop}[単拡大からの埋込みの延長]
体の埋込み $\sigma:K\to E$ と $K$ 上代数的な元 $\alpha$ に対し，$p$ をその最小多項式とする．$\sigma$ を延長する埋込み $K(\alpha)\to E$ は，$\sigma(p)$ の $E$ 内の相異なる根と一対一に対応する．（参照：\GenronRef{K-001307}，\GenronRef{K-001308}．）
\end{prop}
\begin{proof}
延長は $\alpha$ の像 $\beta$ を決めれば一意であり，$\sigma(p)(\beta)=0$ が必要である．逆に根 $\beta$ が与えられれば，係数に $\sigma$ を適用して $\beta$ を代入する環準同型の核は $(p)$ を含む．$K[X]/(p)$ が体であり単位元の像が $1$ なので，誘導される準同型は単射である．単拡大の表示によりこれが延長を与える．
\end{proof}
% genron:end K-001408
% genron:begin K-001409
% genron:use K-000001
% genron:use K-000002
% genron:use K-000003
% genron:use K-000004
% genron:use K-000005
% genron:use K-000007
% genron:use K-000017
% genron:use K-000024
% genron:use K-001406
% genron:use K-001408
% genron:use K-001237
% genron:use K-001307
\begin{prop}[分解体の一意性]
$f\in K[X]$ の任意の二つの分解体は $K$ 上同型である．（参照：\GenronRef{K-001406}．）
\end{prop}
\begin{proof}
二つの分解体を $L,E$ とし，$L$ の根を順次添加する．$K$ の恒等写像から始め，既に定めた部分体の埋込みを次の根に延長する．その根の最小多項式は $f$ を割り，係数に埋込みを適用した多項式も $f$ を割る．$E$ では $f$ が一次式に分解するため根を持つ．前命題により延長できる．有限回の操作で $L\to E$ が得られる．$f$ の分解式の全係数にこの写像を適用すれば，$E$ の全根が像に属することが分かる．$E$ はこれらの根で生成されるので全射である．
\end{proof}
% genron:end K-001409
% genron:begin K-001410
% genron:use K-000001
% genron:use K-000002
% genron:use K-000003
% genron:use K-000004
% genron:use K-000005
% genron:use K-000007
% genron:use K-000017
% genron:use K-000024
% genron:use K-001190
% genron:use K-001306
% genron:use K-001312
\begin{defn}[代数閉体と代数閉包]
体 $E$ が代数閉体であるとは，$E[X]$ の正次数多項式がすべて $E$ 内に根を持つことをいう．$K$ の代数的拡大で代数閉体となるものを $K$ の代数閉包という．（参照：\GenronRef{K-001190}，\GenronRef{K-001306}．）
\end{defn}
% genron:end K-001410
% genron:begin K-001411
% genron:use K-000001
% genron:use K-000002
% genron:use K-000003
% genron:use K-000004
% genron:use K-000005
% genron:use K-000007
% genron:use K-000017
% genron:use K-000024
% genron:use K-001407
% genron:use K-001410
% genron:use K-000129
% genron:use K-000124
% genron:use K-001196
% genron:use K-001311
% genron:use K-001312
% genron:use K-000060
\begin{prop}[代数閉包の存在]
任意の体 $K$ は代数閉包を持つ．（参照：\GenronRef{K-001410}．）
\end{prop}
\begin{proof}
まず体 $F$ に対し，その正次数モニック多項式の集合を整列する．各多項式を順に取り，現在の体上の既約因子の剰余体によって一つの根を添加する．極限段階ではそれまでの体の合併を取る．各剰余体には現在の体が標準的に埋め込まれるので，体を同型なコピーで置き換えて文字どおりの包含を保てる（新しい元には段階の添字を付ければよい）．整列集合上の再帰により，すべての $F$ 上の多項式に根を与える代数的拡大 $F'$ が得られる．各段階で代数性が保たれることは，代数的な係数が有限個ならばそれらで生成される体が有限次となり，さらに代数元を加えても有限次であることから従う．極限段階の各元は以前の段階に属する．
$K_0=K$ とし，上記の操作を繰り返して $K_{n+1}=K_n'$ を取る．この反復は選択を用いるが，根添加の構成は既約因子の選択及び整列を固定すれば再帰的に行える．合併 $E=\bigcup_{n\geq0}K_n$ は体であり $K$ 上代数的である．$E$ 上の正次数多項式の係数は有限個なので一つの $K_n$ に属し，$K_{n+1}$ 内に根を持つ．従って $E$ は代数閉包である．
\end{proof}
% genron:end K-001411
% genron:begin K-001412
% genron:use K-000001
% genron:use K-000002
% genron:use K-000003
% genron:use K-000004
% genron:use K-000005
% genron:use K-000007
% genron:use K-000017
% genron:use K-000024
% genron:use K-001408
% genron:use K-001410
% genron:use K-000131
% genron:use K-001306
% genron:use K-001307
% genron:use K-001190
\begin{prop}[代数閉包への埋込みの延長]
$L/K$ を代数拡大，$E$ を代数閉体，$\sigma:K\to E$ を埋込みとする．$\sigma$ は $L\to E$ に延長できる．特に代数閉包は $K$ 上同型を除いて一意である．（参照：\GenronRef{K-001190}，\GenronRef{K-001306}，\GenronRef{K-001410}．）
\end{prop}
\begin{proof}
$K\subseteq F\subseteq L$ なる中間体と $\sigma$ を延長する埋込み $F\to E$ の組を，延長関係によって順序付ける．鎖の上界は中間体の合併と互いに一致する埋込みの合併である．Zornの補題により極大な組が存在する．$F\ne L$ なら $\alpha\in L\setminus F$ を取り，その最小多項式の像が $E$ に根を持つから単拡大への延長が可能となり矛盾する．従って $F=L$ である．二つの代数閉包の間の埋込みの像は代数閉体であり，第二の代数閉包はその像の上で代数的である．代数閉体の既約多項式は一次なのでその代数拡大は自明であり，埋込みは全射となる．
\end{proof}
% genron:end K-001412
